% ---------------------------------------------------------------------
%  03 · Typographie
%  Schibsted Grotesk pour les titres · Onest pour le texte et les chiffres.
%  Polices embarquées dans fonts/ (licence OFL) : compiler depuis la
%  racine docs/ pour que les chemins se résolvent.
% ---------------------------------------------------------------------

\defaultfontfeatures{Ligatures=TeX}

% Texte courant — Onest (pas d'italique dessinée : penché synthétique)
\setmainfont{Onest}[
  Path           = fonts/onest/,
  Extension      = .ttf,
  UprightFont    = *-Regular,
  BoldFont       = *-SemiBold,
  ItalicFont     = *-Regular,
  BoldItalicFont = *-SemiBold,
  ItalicFeatures     = {FakeSlant=0.16},
  BoldItalicFeatures = {FakeSlant=0.16},
]
\setsansfont{Onest}[
  Path           = fonts/onest/,
  Extension      = .ttf,
  UprightFont    = *-Regular,
  BoldFont       = *-SemiBold,
  ItalicFont     = *-Regular,
  BoldItalicFont = *-SemiBold,
  ItalicFeatures     = {FakeSlant=0.16},
  BoldItalicFeatures = {FakeSlant=0.16},
]
\IfFontExistsTF{DejaVu Sans Mono}
  {\setmonofont{DejaVu Sans Mono}[Scale=MatchLowercase]}
  {\setmonofont{Latin Modern Mono}}

% Titres — Schibsted Grotesk, approche resserrée comme sur les maquettes
\newfontfamily\policeAffiche{SchibstedGrotesk}[
  Path = fonts/schibsted-grotesk/, Extension = .ttf,
  UprightFont = *-Black,     ItalicFont = *-BlackItalic,
  BoldFont    = *-Black,     BoldItalicFont = *-BlackItalic,
  LetterSpace = -3,
]
\newfontfamily\policeTitre{SchibstedGrotesk}[
  Path = fonts/schibsted-grotesk/, Extension = .ttf,
  UprightFont = *-ExtraBold, ItalicFont = *-ExtraBoldItalic,
  BoldFont    = *-Black,     BoldItalicFont = *-BlackItalic,
  LetterSpace = -2,
]
\newfontfamily\policeIntertitre{SchibstedGrotesk}[
  Path = fonts/schibsted-grotesk/, Extension = .ttf,
  UprightFont = *-Bold,      ItalicFont = *-BoldItalic,
  BoldFont    = *-ExtraBold, BoldItalicFont = *-ExtraBoldItalic,
  LetterSpace = -1,
]

% Libellés capitales espacées (« CARTE DU JEUDI 5 NOVEMBRE »)
\newfontfamily\policeLabel{Onest}[
  Path = fonts/onest/, Extension = .ttf,
  UprightFont = *-Medium, BoldFont = *-Bold,
  LetterSpace = 9,
]
% Onest moyen / demi-gras, chiffres tabulaires pour les colonnes
\newfontfamily\policeMoyenne{Onest}[
  Path = fonts/onest/, Extension = .ttf,
  UprightFont = *-Medium, BoldFont = *-Bold,
]
\newfontfamily\policeChiffres{Onest}[
  Path = fonts/onest/, Extension = .ttf,
  UprightFont = *-Medium, BoldFont = *-Bold,
  Numbers = Monospaced,
]

% Raccourcis
% \textlabel[couleur]{texte} : la couleur reste hors de \MakeUppercase
\newcommand\textlabel[2][]{{\if\relax\detokenize{#1}\relax\else\color{#1}\fi
  \policeLabel\footnotesize\MakeUppercase{#2}}}
\newcommand\textmoyen[1]{{\policeMoyenne #1}}
\newcommand\textsecondaire[1]{{\color{texte2}#1}}
\newcommand\deuxchiffres[1]{\ifnum#1<10 0\fi\number#1}
% Numéro de chapitre affiché : 01, 02… puis A, B… en annexe
\newcommand\numchapitre{\deuxchiffres{\value{chapter}}}
\newcommand\numlabel[1]{\IfInteger{#1}{\deuxchiffres{#1}}{#1}}

% Paragraphes
\linespread{1.18}
\setlength\parindent{0pt}
\setlength\parskip{0.55em plus 0.15em minus 0.1em}
\setlength\emergencystretch{2em}
